\subsection*{语法速览}

\begin{table}[htbp]\centering\small
\begin{tabularx}{\linewidth}{@{}lX@{}}
\toprule
\textbf{语法} & \textbf{含义} \\
\midrule
\texttt{mc = mma.fill 0.0} & 把累加器分块初始化为零 \\
\texttt{mma.fill.f16 0.0f} / \texttt{mma.fill.f32 0.0f} & 显式指定精度的累加器 \\
\texttt{mma.load.swiz\textless{}N\textgreater{}} src & 按重排模式加载 \\
\texttt{mma.row.row mc, ma, mb} & C += A × B（均为行主序） \\
\texttt{mma.row.col mc, ma, mb} & C += A × B（A 行主序，B 列主序） \\
\texttt{mma.row.row.sp mc, ma, mb, me} & 带元数据操作数的稀疏 MMA \\
\texttt{mma.row.row.scale mc, ma, mb, sc\_a, sc\_b} & 融合的 MMA 与逐分块反量化 \\
\texttt{mma.scale mc, sc\_a, sc\_b} & 独立的 MMA 后置缩放 \\
\texttt{mma.store.transp mc, dst} & 转置后写出累加器 \\
\texttt{cdiv(a, b)} & 向上取整除法 \\
\texttt{\_\_co\_\_ void fn(...)} & 就地写出结果的核函数 \\
\bottomrule
\end{tabularx}
\end{table}

阅读修饰符时，应同时追踪数据路径。交错装入必须匹配复制产生的布局，稀疏操作需要同一操作数分块的元数据，缩放值也必须对应正在累加的分块。语句写得简短，并不意味着可以忽略这些关系。
